% theory_network_storage_flexibility.tex -- 多时段网络、储能与灵活资源理论

\section{多时段 AC/DC 网络、储能与灵活资源理论}\label{sec:ts-th-network}

\begin{theorynote}
本章为通用理论，讨论时序调度中的 AC/DC 网络近似、储能状态方程、可再生消纳、需求响应、
VSC/DC--DC/多端路由器和移动储能。重点是跨时段守恒、符号和近似误差，而非重复潮流手册的
单时刻非线性方程。与当前实现的对应见章末；未实现模型就地标明。
\end{theorynote}

\subsection{多时段耦合问题的一般形式}
令 $x_t$ 为时段 $t$ 的网络代数状态，$s_t$ 为跨时段状态（SOC、机组状态、温度等），
$v_t$ 为控制量。确定性时序优化可写成
\begin{align}
\min_{x,s,v}\quad&\sum_{t=0}^{T-1}\ell_t(x_t,s_t,v_t)+\Phi(s_T),\label{eq:ts-ocp-obj}\\
\text{s.t.}\quad&g_t(x_t,s_t,v_t)=0,\quad h_t(x_t,s_t,v_t)\le0,\label{eq:ts-ocp-net}\\
&s_{t+1}=F_t(s_t,v_t),\quad s_0=\bar s_0.\label{eq:ts-ocp-state}
\end{align}
若 $g_t$ 为 AC 潮流，则问题是含整数变量的非凸 MINLP；以 DC 潮流、LinDistFlow 或运输网络替代，
可得到 MILP/SOCP。模型层级必须随结果报告，不得仅写“考虑网络约束”。

\subsection{AC 网络：DC 近似、AC 校核和偏差来源}
对支路 $ij$ 的 AC 有功
\begin{equation}
P_{ij}=V_i^2g_{ij}-V_iV_j[g_{ij}\cos\theta_{ij}+b_{ij}\sin\theta_{ij}],\label{eq:ts-acbranch}
\end{equation}
在 $V\approx1$、$r/x\approx0$、$|\theta_{ij}|\ll1$ 下得到
$P_{ij}\approx(\theta_i-\theta_j-\phi_{ij})/x_{ij}$。被删除的主要项为电阻损耗
$g_{ij}\theta_{ij}^2$、电压幅值偏差和高阶正弦项。因此 DC-SCUC 可筛查有功拥塞，
但不能证明 AC 电压、无功、损耗或换流器电流可行。

\begin{proposition}[UC--PF 校核不是等价证明]\label{prop:ts-ucpf}
设某调度满足无损 DC 平衡和支路热限。除非另有条件保证 AC 潮流存在且全部电压、无功和视在功率约束满足，
该调度不能推出 AC 可行。逐时 PF 收敛只证明给定注入存在一个潮流解，也不证明 OPF 约束或调度最优性。
\end{proposition}
\begin{proof}
取高 $R/X$ 两母线系统。DC 模型忽略 $r$ 和电压幅值，任意低于有功热限的负荷均可能满足；
但恒功率负荷增大时 AC 方程可在鞍结点后无解，或在有解时违反电压下限。故 DC 可行不蕴含 AC 可行。
PF 仅解等式，若未把运行不等式加入问题，收敛也不蕴含这些不等式成立。
\end{proof}

可靠流水线应对每时段报告
\begin{equation}
\epsilon_t^V=\max_i|V_{i,t}^{opf}-V_{i,t}^{pf}|,\quad
\epsilon_t^\theta=\max_i|\tilde\theta_{i,t}^{opf}-\tilde\theta_{i,t}^{pf}|,
\quad\epsilon_t^L=P_{loss,t}^{pf}-P_{loss,t}^{opf},\label{eq:ts-crossval}
\end{equation}
其中角度先相对同一参考归一。还应单独检查 $V$ 界、支路 MVA、换流器容量和功率平衡残差。

\subsection{DC 网络与 AC/DC 耦合}
电阻性 DC 网的物理稳态为
\begin{equation}
P_{i,t}^{dc}=V_{i,t}^{dc}\sum_jG_{ij}(V_{i,t}^{dc}-V_{j,t}^{dc}),\label{eq:ts-dcphysical}
\end{equation}
这是二次非线性方程。常见线性近似包括：
\begin{enumerate}
\item 在 $V^{dc}=V^0+\Delta V$ 附近一阶线性化；
\item 无损运输模型 $f_{ij,t}^{dc}\in[-\bar f_{ij},\bar f_{ij}]$ 与节点守恒；
\item 铜板或逐节点自平衡模型，不表达支路传输。
\end{enumerate}
三者物理含义不同。运输模型为
\begin{align}
\sum_{k\in\mathcal C_i}p_{k,t}^{dc}+p_{i,t}^{src}-d_{i,t}^{dc}
&=\sum_{j:(j,i)}f_{ji,t}^{dc}-\sum_{j:(i,j)}f_{ij,t}^{dc},\label{eq:ts-dctransport}\\
|f_{ij,t}^{dc}|&\le\bar f_{ij}^{dc}.\label{eq:ts-dclimit}
\end{align}
它保留连通性和容量，但删除电压、欧姆定律与损耗。

对效率 $\eta_c\in(0,1]$ 的 VSC，若 $p^{ac}>0$ 表示向 AC 注入，则可用方向变量拆分：
\begin{align}
p^{ac}&=p^+-p^-,\quad p^+,p^-\ge0,\\
p^{dc}&=-p^+/\eta_c+\eta_c p^-,\label{eq:ts-vsc-eff}
\end{align}
并用二元方向或凸包约束避免同时双向流。若只用一个固定效率关系，必须明确正负方向的定义。
DC--DC 端口和多端路由器可统一为端口流：
\begin{equation}
\sum_{j\in\mathcal P_r}\eta_{rj}p_{rj,t}^{in}
=\sum_{j\in\mathcal P_r}p_{rj,t}^{out},\qquad
0\le p_{rj,t}^{in/out}\le\bar p_{rj}.\label{eq:ts-router}
\end{equation}

\subsection{储能的能量守恒与符号}
令 $p_t=p_t^{dis}-p_t^{ch}$，$p_t^{ch},p_t^{dis}\ge0$，$E$ 为额定能量，
$e_t\in[\underline e,\overline e]$ 为
\textbf{时段末}能量，$\rho_t=(1-\sigma)^{\Delta t_t}$ 为自放电保持率，则
\begin{equation}
e_t=\rho_te_{t-1}+\eta_c p_t^{ch}\Delta t_t
-\frac{p_t^{dis}\Delta t_t}{\eta_d}.\label{eq:ts-storage-energy}
\end{equation}
若状态用 $SOC_t=e_t/E$，各功率项必须除以 $E$。禁止同时充放电的精确线性化为
\begin{equation}
0\le p_t^{dis}\le\bar P^{dis}z_t,qquad
0\le p_t^{ch}\le\bar P^{ch}(1-z_t),\quad z_t\in\{0,1\}.\label{eq:ts-storage-mode}
\end{equation}

\begin{lemma}[有损储能的同时充放电消除]\label{lem:ts-nosim}
若边际供电成本严格为正、$\eta_c\eta_d<1$、无必须消纳或最小出力过剩，且目标不奖励吞吐量，
则存在最优解满足 $p_t^{ch}p_t^{dis}=0$。
\end{lemma}
\begin{proof}
若二者均正，取足够小的 $\delta>0$ 同时减少 $p_t^{ch}$ 与 $p_t^{dis}$，并调整净注入使状态变化不变。
由于往返损耗为正，维持相同末态所需的外部供电下降，严格降低正边际成本，故原解不可能唯一最优。
若存在负价、弃能惩罚或最小出力过剩，该论证失效，必须保留式~\eqref{eq:ts-storage-mode}。
\end{proof}

周期边界 $e_{T-1}=e_{-1}$ 消除终端能量透支，适合可重复代表日；但它也禁止跨日、跨周和季节性净搬移。
更一般的终端价值方法为
\begin{equation}
\Phi(e_{T-1})=-\lambda_Ee_{T-1}\quad\text{或}\quad
\underline e_T\le e_{T-1}\le\overline e_T,\label{eq:ts-terminal}
\end{equation}
其中 $\lambda_E$ 来自更长时域的水值/储能价值函数。

吞吐量近似循环老化可写为
\begin{equation}
E_t^{through}=\Delta t_t(p_t^{ch}+p_t^{dis}),\qquad
C_t^{deg}=c^{through}E_t^{through},\quad
c^{through}=\frac{C^{replace}}{2N^{cycle}E}.\label{eq:ts-degradation}
\end{equation}
这是等效满循环线性近似，不包含 DoD、温度、C-rate 与雨流计数路径依赖。

\subsection{可再生、需求响应与外部电网}
可再生出力满足
\begin{equation}
0\le p_{rt}^{res}\le\bar p_{rt}^{avail},\qquad
p_{rt}^{curt}=\bar p_{rt}^{avail}-p_{rt}^{res}.\label{eq:ts-renewable}
\end{equation}
must-take 模式等价于固定 $p^{curt}=0$，在过剩时可能使模型不可行。外部电网交换
$p_t^{ext}\in[-\bar P^{export},\bar P^{import}]$ 应用分时价格结算；若使用任意大 Big-M，
其大小必须由变电站容量或网络切集给出，否则数值条件和物理边界均不可审计。

需求响应可拆为增、减负荷变量：
\begin{align}
d_{ft}^{served}&=d_{ft}^0+p_{ft}^{up}-p_{ft}^{down},\label{eq:ts-dr}\\
0\le p_{ft}^{up}&\le\bar p_{ft}^{up},\quad0\le p_{ft}^{down}\le\bar p_{ft}^{down},\\
\sum_t(p_{ft}^{up}-p_{ft}^{down})\Delta t_t&=0.\label{eq:ts-dr-energy}
\end{align}
式~\eqref{eq:ts-dr-energy} 表示可移峰负荷；若删除则是可中断/增负荷服务。二者不能共用同一业务解释。

\subsection{移动储能与聚合资源}
移动储能位置变量 $a_{ibt}\in\{0,1\}$ 满足 $\sum_ba_{ibt}+w_{it}=1$，
$w_{it}$ 为在途状态。母线注入受位置门控
\begin{equation}
-\bar P^{ch}a_{ibt}\le p_{ibt}\le\bar P^{dis}a_{ibt},\qquad
e_{it}=F(e_{i,t-1},\sum_bp_{ibt})-E_i^{travel}w_{it}.\label{eq:ts-mobile}
\end{equation}
完整模型还需道路可达、旅行时间、最短停留和唯一位置守恒。VPP 的聚合功率/能量包络是内部资产的外近似；
若没有逐资产可实现性约束，聚合可行并不必然存在内部可行分解。

\subsection{六步外部网络交叉验证}
TS-PF-6STEP 使用 10 MVA、12.47 kV 两母线系统，支路 $z=0.012+j0.032$ pu，
基准负荷 $1.2+j0.45$ MVA，倍率为 $(0.65,0.8,1,1.15,1.3,0.9)$。
六步 HySim 负荷母线电压依次为
\begin{equation*}
(0.998122169,\ 0.997687158,\ 0.997106165,\ 0.996669684,\ 0.996232570,\ 0.997396802)\ \mathrm{pu}.
\end{equation*}
相同物理量经 OpenDSS 独立求解，最大误差为
$6.59\times10^{-10}$ pu、$4.51\times10^{-8}$ 度、
$6.64\times10^{-7}$ MW 和 $2.59\times10^{-7}$ Mvar；
GridLAB-D 最大误差为 $5.85\times10^{-7}$ pu、$4.17\times10^{-7}$ 度、
$4.81\times10^{-6}$ MW 和 $4.99\times10^{-7}$ Mvar，6/6 通过。

\subsection{与实现的对应}
\begin{center}\footnotesize
\begin{longtable}{@{}L{6.3cm}L{8.9cm}@{}}
\toprule 理论条目 & 实现状态\\\midrule\endfirsthead
\toprule 理论条目 & 实现状态\\\midrule\endhead
UC 阶段 DC 相角/支路约束 & \implfull{src/time\_series/time\_series\_pf.cpp:build\_uc\_milp}\\
逐时 AC-OPF 与 PF 校核，式~\eqref{eq:ts-crossval} &
\implfull{time\_series\_pf.cpp:solve\_time\_series\_pf}\\
物理 DC 电压方程，式~\eqref{eq:ts-dcphysical} & \implnone\\
显式 DC 运输流与热限，式~\eqref{eq:ts-dctransport}--\eqref{eq:ts-dclimit} &
\implfull{build\_uc\_milp；需 enable\_dc\_branch\_flows=true}\\
VSC/DC--DC 效率耦合 & \implpart{有方向变量与效率关系；不是 DC 电压物理模型}\\
储能有损 SOC、自放电、充放电互斥和周期边界 &
\implfull{build\_uc\_milp；tests/test\_uc\_storage\_efficiency.cpp}\\
吞吐量退化费和日循环上限，式~\eqref{eq:ts-degradation} &
\implfull{build\_uc\_milp；enable\_storage\_degradation}\\
雨流、温度、DoD/C-rate 电化学退化 & \implnone\\
可再生削减、外部电网、DR 能量守恒 &
\implfull{build\_uc\_milp；对应 opt-in 选项}\\
移动储能外生与联合选址状态 & \implpart{两种路径均实现；道路网旅行时间不是本模块内生模型}\\
VPP 聚合功率/能量包络 & \implpart{实现聚合包络；未证明内部资产可分解}\\
OpenDSS/GridLAB-D 六步快照矩阵 &
\implfull{tools/time\_series\_validation/run\_cross\_engine\_matrix.py}\\
\bottomrule
\end{longtable}
\end{center}

\subsection{参考文献}
{\renewcommand{\section}[2]{}\begin{thebibliography}{9}\small
\bibitem{ts:kersting} W.~H. Kersting, \emph{Distribution System Modeling and Analysis}, 4th ed., CRC Press, 2017.
\bibitem{ts:baran} M.~E. Baran and F.~F. Wu, Optimal sizing of capacitors placed on a radial distribution system,
\emph{IEEE Transactions on Power Delivery}, 4(1), 1989.
\bibitem{ts:storage} A.~Castillo and D.~F. Gayme, Grid-scale energy storage applications in renewable energy integration,
\emph{Energy Conversion and Management}, 87, 2014.
\end{thebibliography}}
